% ======================================================================
% SEÇÃO 2 — FUNDAMENTAÇÃO E TRABALHOS RELACIONADOS
% ======================================================================
\section{Fundamenta\c{c}\~ao e trabalhos relacionados}
\label{sec:fund}

\subsection{Especia\c{c}\~ao de merc\'urio e instrumenta\c{c}\~ao anal\'itica}

A determina\c{c}\~ao de esp\'ecies de merc\'urio em concentra\c{c}\~oes tra\c{c}o
exige etapas de deriva\c{c}\~ao qu\'imica que convertam as esp\'ecies em formas
vol\'ateis e est\'aveis, adequadas \`a separa\c{c}\~ao e \`a detec\c{c}\~ao.
Bloom \cite{bloom1989} e Liang e colaboradores \cite{liang1994} estabeleceram os
protocolos de etila\c{c}\~ao em fase aquosa com NaBEt$_4$ e de separa\c{c}\~ao
por cromatografia gasosa com detec\c{c}\~ao por fluoresc\^encia at\^omica em vapor
frio, que s\~ao a base qu\'imica do processo aqui automatizado. Revis\~oes como a
de Leermakers e colaboradores \cite{leermakers2005} evidenciam que a
repetibilidade das condi\c{c}\~oes de prepara\c{c}\~ao --- em particular da rampa
t\'ermica --- \'e um fator cr\'itico de qualidade dos m\'etodos de
especia\c{c}\~ao. No plano instrumental, equipamentos comerciais dedicados, como
o analisador RA-915 da Lumex \cite{lumex}, integram a atomiza\c{c}\~ao e a
detec\c{c}\~ao, mas tipicamente n\~ao exp\~oem ao usu\'ario o controle fino do
sequenciamento da prepara\c{c}\~ao que antecede a medida --- refor\c{c}ando a
relev\^ancia de um sistema de automa\c{c}\~ao dedicado a essa etapa.

\subsection{Controle de processos e rampas de temperatura}

Em processos cuja vari\'avel controlada deve seguir uma trajet\'oria no tempo,
como a rampa linear de temperatura do Tubo U, o controlador
proporcional-integral-derivativo (PID) \'e a solu\c{c}\~ao mais difundida
\cite{astrom2006,ogata2010}. Em sua forma cont\'inua ideal, a sa\'ida \'e dada
por
\begin{equation}
  u(t) = K_p \left[ e(t) + \frac{1}{T_i}\int_0^t e(\tau)\,d\tau
         + T_d \frac{d\,e(t)}{dt} \right],
  \label{eq:pid}
\end{equation}
em que $K_p$ \'e o ganho proporcional, $T_i$ o tempo integral e $T_d$ o tempo
derivativo. Implementa\c{c}\~oes digitais devem tratar a satura\c{c}\~ao da sa\'ida
(anti-windup) e aplicar a a\c{c}\~ao derivativa sobre a vari\'avel de processo,
evitando o \emph{derivative kick} \cite{astrom2006}. Quando a medi\c{c}\~ao n\~ao
\'e confi\'avel em parte da faixa --- como abaixo do piso de leitura do termopar
---, uma estrat\'egia de raz\~ao de taxas em malha aberta pode ser combinada ao
PID, calculando a vari\'avel manipulada proporcionalmente \`a rela\c{c}\~ao entre
a taxa de aquecimento desejada e a taxa do sistema.

O controle sequencial de processos \'e frequentemente modelado por
\emph{m\'aquinas de estado finitas} (FSM), que associam a cada estado um conjunto
de sa\'idas ativas (matriz de acionamento) e tratam de forma sistem\'atica
eventos de parada e emerg\^encia \cite{wagner2006}.

\subsection{Plataformas embarcadas e comunica\c{c}\~ao}

A literatura e a pr\'atica de instrumenta\c{c}\~ao anal\'itica adotam, com
crescente frequ\^encia, plataformas abertas de baixo custo: microcontroladores
como o Arduino \cite{arduino} e computadores de placa \'unica como o Raspberry Pi
\cite{raspberrypi}. A separa\c{c}\~ao entre uma camada de supervis\~ao de alto
n\'ivel e um m\'odulo de aquisi\c{c}\~ao em tempo real melhora a confiabilidade e
a manutenibilidade. A medi\c{c}\~ao de temperatura por termopares tipo K \'e
comumente digitalizada por circuitos integrados com compensa\c{c}\~ao de
jun\c{c}\~ao fria, como o MAX6675, que entrega 12~bits com resolu\c{c}\~ao de
$0{,}25\,\celsius$ via SPI \cite{max6675,nist_its90}.

A comunica\c{c}\~ao entre plataformas \'e realizada com pacotes JSON
\cite{bray2017} em enlace serial, e a IHM Web consome uma API REST e telemetria
em tempo real via WebSocket \cite{fette2011}. Frameworks modernos como FastAPI
\cite{fastapi} e React \cite{react} permitem construir supervis\'orios acess\'iveis
por navegador, em contraste com plataformas propriet\'arias monol\'iticas
(frequentemente baseadas em LabVIEW) que acoplam controle, interface e
aquisi\c{c}\~ao em um \'unico ambiente, com custo elevado de licenciamento e
evolu\c{c}\~ao restrita.

\subsection{S\'intese e lacuna}

A revis\~ao indica que n\~ao h\'a, no dom\'inio considerado, uma solu\c{c}\~ao
aberta e de baixo custo que integre, de forma coesa e segura, o sequenciamento da
prepara\c{c}\~ao de amostras, o controle PID da rampa do Tubo U e uma IHM Web com
persist\^encia de par\^ametros. Este artigo busca preencher essa lacuna, com
\^enfase na rastreabilidade entre requisitos, implementa\c{c}\~ao e valida\c{c}\~ao
por testes.
